%! Sinusoidal Function Context and Data Modeling — Unit 6, Lesson 6.8 — Lesson Plan
\documentclass[10pt]{article}
\usepackage{precalculus-article}
\usepackage{precalculus-boxes}
\usepackage{graphicx}
\graphicspath{{images/}}
\newcommand{\TallMath}[1]{$\displaystyle #1\rule[-1.4em]{0pt}{3.2em}$}
\newcommand{\CourseName}{Precalculus}
\newcommand{\MeetingLength}{55 minutes}
\newcommand{\UnitNumberName}{Unit 6: Trigonometric Functions \quad}
\newcommand{\LessonNumberName}{Lesson 6.8: Sinusoidal Function Context and Data Modeling}

\begin{document}

%----------------------------------------------------------------------
% Title Block
%----------------------------------------------------------------------
\begin{center}
  {\Large\bfseries \CourseName} \\[4pt]
  {\small\bfseries \UnitNumberName \LessonNumberName \quad (2--3 instructional periods, \MeetingLength\ each)}
\end{center}

%----------------------------------------------------------------------
% Primary Objective
%----------------------------------------------------------------------
\begin{tcolorbox}[colback=sky, colframe=navy, boxrule=0.9pt, arc=2mm,
    left=2mm, right=2mm, top=1.5mm, bottom=1.5mm]
  \textbf{Primary Objective:} Students will construct sinusoidal function models
  of periodic phenomena by extracting amplitude, vertical shift, period, and
  phase shift from contextual data, and will use those models to predict values
  within a contextual domain.
  \textit{(Big Idea: Functions; Skills 1.C, 3.C)}
\end{tcolorbox}

\vspace{0.08in}

%----------------------------------------------------------------------
% Priority Ideas & Skills
%----------------------------------------------------------------------
\begin{skillbox}[Priority Ideas \& Skills]{goldbox}
\begin{tabularx}{\linewidth}{@{} p{0.46\linewidth} X @{}}
\textbf{AP Skills} & \textbf{Key Understandings (EKs)} \\
\midrule
\textbf{1.C} \textit{Construct new functions} ---
  Build a sinusoidal model from a data set by estimating key parameter values
  or by performing a sinusoidal regression with technology.
&
\textbf{EK 3.7.A.1}: The smallest input-value interval between consecutive
  maxima (or consecutive minima) in the data estimates the period. \\[4pt]

\textbf{3.C} \textit{Support conclusions with data} ---
  Evaluate the reasonableness of a sinusoidal model by comparing predicted
  values to actual data, and clarify the contextual domain over which the
  model is useful.
&
\textbf{EK 3.7.A.2}: Amplitude $= \tfrac{\max - \min}{2}$; vertical
  shift/midline $= \tfrac{\max + \min}{2}$. \\[4pt]

&
\textbf{EK 3.7.A.3}: A known input-output pair is compared to the model to
  estimate the phase shift. \\[4pt]

&
\textbf{EK 3.7.A.4}: Technology can be used to perform sinusoidal regression
  and to estimate parameters from data. \\[4pt]

&
\textbf{EK 3.7.A.5}: Sinusoidal models are valid only over a contextual domain
  and are used to predict output values from input values. \\
\end{tabularx}
\end{skillbox}

\vspace{0.08in}

%----------------------------------------------------------------------
% Vocabulary
%----------------------------------------------------------------------
\begin{skillbox}[{Vocabulary, Concepts, \& Theorems}]{greenbox}
\begin{tabularx}{\linewidth}{@{} p{0.26\linewidth} X @{}}
\textbf{Term} & \textbf{Definition / Note} \\
\midrule
sinusoidal model &
  A function of the form $f(x) = A\sin(B(x - C)) + D$ (or cosine form)
  used to represent a periodic real-world quantity. \\[4pt]
amplitude &
  Half the range of output values: $A = \tfrac{\max - \min}{2}$. Measures
  how far the function deviates from its midline. \\[4pt]
vertical shift (midline) &
  $D = \tfrac{\max + \min}{2}$; the horizontal axis about which the
  sinusoidal function oscillates. \\[4pt]
period &
  The input-value distance between consecutive maxima (or minima) in
  the data; equals $\tfrac{2\pi}{B}$ in the model. \\[4pt]
phase shift &
  $C$; determined by comparing a known data point (e.g., a maximum) to
  the standard sine/cosine maximum location. \\[4pt]
contextual domain &
  The restricted set of input values over which the model is meaningful
  given the real-world situation. \\
\end{tabularx}
\end{skillbox}

\vspace{0.08in}

%----------------------------------------------------------------------
% Activate Prior Knowledge
%----------------------------------------------------------------------
\begin{skillbox}[Activate Prior Knowledge \& Spiral Review]{sky}
\begin{tabularx}{\linewidth}{@{}X@{}}
\textbf{Spiral topics activated during Warm-Up (first 8--10 minutes):}
\begin{itemize}[leftmargin=1.4em, itemsep=2pt]
  \item \textbf{Reading sinusoidal graphs (Lessons 3.5--3.6):} Identify amplitude,
    period, and midline from a graph --- connects directly to extracting the same
    quantities from data tables.
  \item \textbf{Computing $(\max + \min)/2$ and $(\max - \min)/2$ from a data set:}
    A prerequisite numerical skill targeted directly in Warm-Up Problem 2.
  \item \textbf{Function transformation parameters $A$, $B$, $C$, $D$:}
    Students should recall how each parameter shifts or scales the parent
    $\sin$ or $\cos$ function before applying those ideas to real data.
\end{itemize}
\end{tabularx}
\end{skillbox}

\vspace{0.08in}

%----------------------------------------------------------------------
% Hook
%----------------------------------------------------------------------
\begin{skillbox}[Hook --- Entry Event]{sky}
\textbf{Scenario:} Display a scatter plot of average monthly high temperature
(${}^\circ$F) for a U.S. city over one calendar year. The data trace a smooth
wave: low in January, peaking in July, dropping again by December.

\textbf{Discussion prompts:}
\begin{enumerate}[leftmargin=1.4em, itemsep=3pt]
  \item \textit{``What kind of function could model this shape? How do you know?''}
  \item \textit{``What information from the table would help you write the equation?''}
  \item \textit{``If we use this model to predict the temperature in month 15,
    should we trust that prediction? Why or why not?''}
\end{enumerate}

\textbf{Key tension to surface:} Students know sinusoidal functions from Lessons
3.5--3.6 but have only worked with equations given directly. Today they must
\emph{extract} parameters from messy real-world data --- a new challenge.

\textbf{Transition:} ``Today we'll learn a systematic four-step process for turning
any periodic data set into a sinusoidal model, then evaluate how well our model fits.''
\end{skillbox}

\vspace{0.08in}

%----------------------------------------------------------------------
% Lesson
%----------------------------------------------------------------------
\begin{skillbox}[Lesson --- Instruction Sequence]{sky}
\begin{multicols}{2}

\textbf{Step 1 --- Estimate Period from Data} (8 min)

Present a 12-month temperature data table. Ask: ``Where are the consecutive
maxima?'' Model how to identify the $x$-values of two consecutive peaks and
compute their difference (EK 3.7.A.1). Then compute $B = \tfrac{2\pi}{k}$.
Verify by checking consecutive minima; both should give the same period.

\columnbreak

\textbf{Step 2 --- Compute Amplitude and Midline} (8 min)

Show the formulas $A = \tfrac{\max - \min}{2}$ and $D = \tfrac{\max + \min}{2}$
(EK 3.7.A.2). Apply to the temperature data. Discuss the midline in context:
``the average between the hottest and coolest month.'' Emphasize that amplitude
is measured \emph{from} the midline, not from zero.

\end{multicols}

\begin{multicols}{2}

\textbf{Step 3 --- Determine Phase Shift} (10 min)

For the sine model $f(x) = A\sin(B(x - C)) + D$: standard sine reaches its
maximum at $B(x - C) = \tfrac{\pi}{2}$, i.e., $x = C + \tfrac{\pi}{2B}$.
Set this equal to the $x$-value of the data's first maximum; solve for $C$
(EK 3.7.A.3). Validate: substitute a second data point to check closeness.
Tip: if a maximum occurs at $x = x_0$, using a cosine model gives $C = x_0$
directly.

\columnbreak

\textbf{Step 4 --- Validate and Predict} (10 min)

Evaluate the model at several data points. Discuss residuals (predicted minus
actual). Use the model to predict a value (e.g., temperature in month 14)
and state the contextual domain (EK 3.7.A.5). Caution: far outside the domain,
real-world factors (climate change, relocation) may break the assumption of
perfect periodicity.

\end{multicols}

\begin{multicols}{2}

\textbf{Step 5 --- Technology: Sinusoidal Regression} (8 min)

Demonstrate entering data in Desmos or a graphing calculator and running a
sinusoidal regression (EK 3.7.A.4). Compare regression parameters to hand
estimates. Discuss why they may differ: the regression minimizes total error
across all points, not just the extremes.

\columnbreak

\textbf{Step 6 --- Contextual Domain Discussion} (5 min)

Ask: ``Should we trust this model to predict temperatures 50 years from now?''
Reinforce EK 3.7.A.5: validity requires the periodic assumption to hold.
Real climate data, population trends, or equipment changes may violate that
assumption outside the observed domain.

\end{multicols}

\smallskip
\textbf{Pacing note:} Steps 1--3 in Period 1; Steps 4--6 + activity in Periods 2--3.
\end{skillbox}

\vspace{0.08in}

%----------------------------------------------------------------------
% Active Monitoring
%----------------------------------------------------------------------
\begin{skillbox}[Active Monitoring]{redbox}
\begin{itemize}[leftmargin=1.4em, itemsep=3pt]
  \item \textbf{Period vs.\ half-period:} Students measure from a maximum to the
    next minimum (half-period). Prompt: \textit{``We need two \emph{peaks} --- maximum
    to maximum, not maximum to minimum.''}
  \item \textbf{Amplitude vs.\ maximum value:} Students write $A = \max$. Prompt:
    \textit{``The amplitude is measured from the midline. What is the midline here?''}
  \item \textbf{Phase shift sign errors:} The form $\sin(B(x - C))$ shifts right by
    $C$ (positive $C$ means shift right). Students flip the sign. Reinforce: substitute
    $x = C$ --- do you get a maximum output?
  \item \textbf{Contextual domain over-extension:} Students apply the model beyond
    the data range without comment. Ask: \textit{``What real-world constraint limits
    meaningful input values for this model?''}
\end{itemize}
\end{skillbox}

\vspace{0.08in}

%----------------------------------------------------------------------
% Group Work & Differentiation
%----------------------------------------------------------------------
\begin{skillbox}[Group Work \& Differentiation]{redbox}
Students work in groups of 3--4 on the tiered activity. Assign tiers or allow
structured self-selection with teacher guidance.

\begin{itemize}[leftmargin=1.4em, itemsep=4pt]
  \item \textbf{Tier R (Reinforce):} Given a complete data table with max, min,
    and the period already labeled, students compute $A$, $D$, and $B$ using
    scaffolded formula templates and fill in the model $f(x) = A\sin(B(x-C))+D$.
    A step-by-step guide and calculator are provided.
  \item \textbf{Tier A (Apply):} Given a tidal data table over 24 hours, students
    identify max, min, period, amplitude, midline, and phase shift; write the full
    model; and use it to predict the tide height at a specified time.
  \item \textbf{Tier E (Extend):} Given temperature data for two cities, students
    build models for both, compare parameters (which city is warmer on average?
    which has greater seasonal swing?), compute residuals, and determine which
    city's climate is more faithfully modeled by a sinusoidal function.
\end{itemize}
\end{skillbox}

\vspace{0.08in}

%----------------------------------------------------------------------
% Individual Work & Assessment
%----------------------------------------------------------------------
\begin{skillbox}[Individual Work \& Assessment]{redbox}
Exit ticket administered independently (final 5 minutes of Period 2 or 3).

\begin{enumerate}[leftmargin=1.4em, itemsep=4pt]
  \item A small 6-value data table is provided. Students identify the maximum,
    minimum, amplitude, and midline.
  \item Students identify the period from consecutive maxima in the table and
    compute $B$.
  \item Students use the given model to predict one output value and state whether
    that input is inside the contextual domain.
\end{enumerate}

\textbf{Data use:} Sort exit tickets by which parameter caused the most errors.
Use the most common error to open Lesson 3.8 with a 60-second targeted recap.
\end{skillbox}

\vspace{0.08in}

%----------------------------------------------------------------------
% Reinforcement & Extension
%----------------------------------------------------------------------
\begin{skillbox}[Reinforcement \& Extension]{goldbox}
\textbf{Homework (Lesson 6.8):} 5--6 problems covering: extracting parameters
from a data table, writing a sinusoidal model, using the model to predict,
evaluating model fit, and 2 AP-style multiple-choice items. An optional extension
asks students to run a sinusoidal regression in Desmos and compare the regression
equation to their hand-built model.

\textbf{Preview of Lesson 3.8 --- The Tangent Function:}
Ask students: \textit{``Sine and cosine model bounded, wave-like phenomena.
What if a periodic quantity shoots off toward infinity rather than staying bounded?
What function might capture that behavior?''} Lesson 3.8 introduces the tangent
function, its vertical asymptotes, and its definition as $\tan\theta = \sin\theta / \cos\theta$.
\end{skillbox}

\end{document}
